\section{项目定位}

\begin{frame}{项目一句话定位}
\begin{block}{研究对象}
在共享内存多核 CPU 上，研究如何把大量有限元单元刚度矩阵高效、正确地累加到全局稀疏刚度矩阵。
\end{block}
\begin{columns}[T,onlytextwidth]
\begin{column}{0.48\textwidth}
\textbf{不是当前主线}
\begin{itemize}
  \item 继续扩展 GPU 算法。
  \item 做完整商业求解器。
  \item 优化线性方程组求解器。
\end{itemize}
\end{column}
\begin{column}{0.48\textwidth}
\textbf{当前主线}
\begin{itemize}
  \item CPU 并行组装算法对比。
  \item 真实工程网格可复现 benchmark。
  \item 结果、图表、文档和短期汇报可追溯。
\end{itemize}
\end{column}
\end{columns}
\source{README.md; docs/requirements/cpu-parallel-stiffness-assembly-design.md}
\end{frame}
\note{
这个项目的核心不是“有限元所有阶段都优化”，而是聚焦 assembly：从每个 element 的局部刚度矩阵 Ke 出发，把它写回 global stiffness matrix K。求解阶段、边界条件完整处理、非线性材料、MPI 和 GPU 新算法都不是当前主线。
}

\begin{frame}{为什么 assembly 值得单独研究}
\begin{itemize}
  \item 有限元里每个单元都能独立计算局部矩阵，看起来天然并行。
  \item 但多个单元会共享节点和自由度，写回全局矩阵时会写到同一位置。
  \item 因此并行 assembly 的难点不是“算 Ke”，而是\term{如何避免或管理写冲突}。
  \item 真实大网格上，全局矩阵很稀疏，不能用 dense matrix 方式处理。
\end{itemize}
\takeaway{并行 assembly 是“局部计算容易、全局写入困难”的问题。}
\source{docs/requirements/cpu-parallel-stiffness-assembly-design.md; ScienceDirect S0045782524003323}
\end{frame}
\note{
从新手角度理解：每个单元像一小块拼图，Ke 是这块拼图内部自由度之间的刚度关系。把所有拼图放回整体结构时，相邻单元共享边界节点，所以它们会同时贡献到同一个全局矩阵条目。并行时这些共享位置会造成 race condition。
}

\begin{frame}{当前唯一主线目录}
\begin{block}{继续开发入口}
\code{parallel\_global\_stiffness\_assembly/cpu\_parallel\_stiffness\_assembly}
\end{block}
\begin{itemize}
  \item 仓库根 README 已明确：当前项目聚焦 CPU 平台并行整体刚度矩阵组装。
  \item GPU 历史内容保留为 legacy 参考，不是当前默认入口。
  \item 短期 beamer 已在 \code{reports/2026-05-14-mentor-next-steps-beamer/}。
  \item 本长期 beamer 固定放在 \code{reports/project-long-term-beamer/}。
\end{itemize}
\takeaway{后续新增知识、稳定结果和进展复盘，合并到长期 deck；日期目录只承载一次汇报。}
\end{frame}
\note{
这一页是为了防止项目材料散掉。当前仓库里有根 README、docs、results、reports、legacy_gpu 等多类资产。长期 deck 的作用就是把这些资产串成一条学习线。
}

\begin{frame}{项目最终要回答的问题}
\begin{enumerate}
  \item 图着色法在多核 CPU 上是否仍然是合适基线？
  \item 直接累加和 atomic 在真实工程网格上是否可接受？
  \item 线程私有缓冲、COO sort-reduce、row-owner 等路线各自代价是什么？
  \item 规则网格和真实 \code{C3D4} 工程网格上的趋势是否一致？
  \item 工程上后续应该优先保留哪类算法继续优化？
\end{enumerate}
\source{docs/requirements/cpu-parallel-stiffness-assembly-design.md}
\end{frame}
\note{
这些问题来自需求文档。注意它们不是孤立算法题，而是研究平台问题：同一输入、同一 CSR、同一 correctness baseline、同一 benchmark schema 下，比较不同 assembly strategy 的性能、内存和正确性。
}

\begin{frame}{当前已经具备的基础}
\begin{columns}[T,onlytextwidth]
\begin{column}{0.52\textwidth}
\textbf{实现基础}
\begin{itemize}
  \item 3D \code{Tet4} / \code{Hex8}。
  \item 每节点 3 个位移 DOF。
  \item CSR 稀疏结构预计算。
  \item \code{AssemblyPlan} 和 scatter plan。
  \item Abaqus \code{.inp} 基础解析。
\end{itemize}
\end{column}
\begin{column}{0.44\textwidth}
\textbf{实验基础}
\begin{itemize}
  \item 串行 correctness baseline。
  \item 多个 CPU 并行 assembler。
  \item CSV / JSON / Markdown 输出。
  \item 结果图表和跨平台 schema。
\end{itemize}
\end{column}
\end{columns}
\source{parallel\_global\_stiffness\_assembly/cpu\_parallel\_stiffness\_assembly/README.md}
\end{frame}
\note{
这里要记住：项目已经不是“从零搭代码骨架”。当前任务更多是把平台变得可解释、可复现、可汇报。长期 deck 的作用之一就是让代码事实、benchmark 事实和汇报话术一致。
}

\begin{frame}{真实工程网格为什么重要}
\begin{itemize}
  \item \code{3d-WindTurbineHub.inp} 是当前核心真实算例。
  \item 规模约为 \metric{228384} nodes、\metric{1113684} elements、\metric{685152} DOFs。
  \item 单元类型为 \code{Tet4} / Abaqus \code{C3D4}。
  \item 真实非结构化网格更能暴露写冲突、稀疏结构构建、内存和缓存行为。
\end{itemize}
\takeaway{小规则网格适合 smoke test；WindTurbineHub 才是工程可用性判断的主要证据。}
\source{results/2026-05-11-symbolic-numeric/symbolic\_numeric\_eval\_report.md}
\end{frame}
\note{
规则 cube 网格好处是可控、快、容易调试；真实 WindTurbineHub 好处是规模和拓扑更接近实际工程。算法在小网格上正确，不代表真实大网格上内存和扩展性可接受。
}

\begin{frame}{短期 beamer 和长期 beamer 的边界}
\vspace{4pt}
\begin{center}
\begin{tabularx}{\textwidth}{>{\bfseries}lXX}
\toprule
类型 & 目的 & 更新原则 \\
\midrule
短期 beamer & 服务一次 mentor/组会/课程汇报 & 围绕日期和听众压缩叙事，可以只选部分证据 \\
长期 beamer & 服务个人掌握完整项目 & 保存稳定概念、结果、来源、进度和待办 \\
结果报告 & 保存实验输出和图表 & 不为演讲改事实，只按数据生成 \\
\bottomrule
\end{tabularx}
\end{center}
\takeaway{长期 deck 不是把短期 deck 堆起来，而是吸收其中经验证的稳定知识。}
\end{frame}
\note{
这个边界很重要。短期汇报可以强调某个 mentor 问题，比如二阶段组装和 Intel 平台优先级；长期 deck 必须把这件事放回项目整体脉络。
}
